% Copyright (c) 2017 by Daniel R. Driver.
% !TEX encoding = UTF-8 Unicode
% !TEX TS-program = XeLaTeX

\documentclass[titlepage]{article}

% This document presumes a file structure and set of inputs that are
% available at: git@github.com:danieldriver/syllabi.git

\newcommand\policy{../policy}
\newcommand\incl{../includes}
\ProvidesFile{variables.tex}[2022/09/11 v2.3 -- Syllabus variables]

\usepackage{xspace} % make manual spaces (like \mycmd\ ) unnecessary
\usepackage{xifthen} % provides \isempty test

% variables for internal use
\newcommand\prof{}
\newcommand\pdegree{}
\newcommand\pphone{}
\newcommand\pemail{}
\newcommand\poffice{}
\newcommand\phours{}
%
\newcommand\ccode{}
\newcommand\ctitle{}
\newcommand\cseries{}
\newcommand\cversion{}
\newcommand\csemester{}
\newcommand\cmeetson{}
\newcommand\cmeetsat{}
\newcommand\cmeetsin{}
\newcommand\cwebsite{}
\newcommand\cdescrip{}
\newcommand\cprereqs{}
\newcommand\edobject{}

% in case of fully online, BTh, or undergraduate courses - https://tex.stackexchange.com/a/5896
\newif\ifonline
\newcommand\Int[2]{\ifonline#1\else#2\fi}
\newif\ifbth
\newcommand\bth[2]{\ifbth#1\else#2\fi}
\newif\ifUG
\newcommand\UG[2]{\ifUG#1\else#2\fi}

% commands for setting variables in the preamble
\newcommand\professor[2][PhD]{
  \renewcommand\pdegree{#1\xspace}
  \renewcommand\prof{#2\xspace}}
\newcommand\phone[1]{
  \renewcommand\pphone{\addfontfeatures{Numbers=Monospaced}#1\xspace}}
\newcommand\email[1]{
  \renewcommand\pemail{\href{mailto:#1}{#1}\xspace}}
\newcommand\officehours[2][AST Library/Teams]{
  \renewcommand\poffice{#1\xspace}
  \renewcommand\phours{\href{https://outlook.office.com/bookwithme/user/b97b147a2e5c43b4a0e200c2cb630b77@astheology.ns.ca?anonymous}{#2}\xspace}%
  }
%
\newcommand\coursecode[2][1.0]{
  \renewcommand\cversion{#1\xspace}
  \renewcommand\ccode{#2\xspace}}
\newcommand\coursetitle[2][]{
  \ifthenelse{\isempty{#1}}%
    {}% do nothing if #1 is empty, else:
    {\renewcommand\cseries{#1\\[1ex]}}
  \renewcommand\ctitle{#2\xspace}}
\newcommand\semester[1]{
  \renewcommand\csemester{#1\xspace}}
\newcommand\meets[3]{
  \newcommand\AM{\textsc{am}}
  \newcommand\PM{\textsc{pm}}
  \renewcommand\cmeetson{#1\xspace}
  \renewcommand\cmeetsat{\Int{From 9:00 \AM}{#2}\xspace}
  \renewcommand\cmeetsin{\Int{\href{https://smu.brightspace.com/d2l/login}{Brightspace}}{#3}\xspace}}
\newcommand\website[2][]{
  \renewcommand\cwebsite{
  \ifthenelse{\isempty{#1}}%
    {\href{https://#2}}% copy url if #1 is empty, else
    {\href{#1}}% use specified url and print #2
  {#2}\xspace}}
\newcommand\cdescription[2][BF 1001; RM 1000 or GTRS 6000]{
  \renewcommand\cprereqs{#1}
  \renewcommand\cdescrip{#2\par}}
\newcommand\objectives[1]{
  \renewcommand\edobject{#1\par}}


\coursecode[2.1.4]{HB 3111}
\coursetitle[Hebrew Bible]{The Book of Genesis}

% Taught as RLGS 3213 (T&I: Genesis) in:
%   - Winter 2009
% Taught similar course BSTH 4403 (TOBITH: Creation and the OT) in:
%   - Fall 2013 (students dropped until it became a DS with Rizalde Dequilla)
% Taught as HB 3111 in:
%   - Winter 2017

\professor{Daniel R. Driver}
\phone{902-425-7051}
\email{ddriver@astheology.ns.ca}
\officehours{Tuesdays, 1:30--3:30 \PM}

\semester{Winter Term 2017}
\meets{Fridays}% \meets{on}{at}{in}
      {9:00--11:30 \AM}
      {Memorial Room, St. Paul’s}
\website{danieldriver.com/courses/}
\cdescription{% copy from the current Academic Calendar
	In the modern period an enormous amount has been written about the
	genesis of Genesis, the history \emph{behind} the first book of the
	Bible. Prior to the critical era, however, perhaps even more has
	been written in commentary upon the Book of Genesis, forming a
	history \emph{before} the text (the history of the book's
	reception). These two approaches to the history of Genesis can seem
	radically different. Indeed, some commentators today still hold that
	the dawn of the historical-critical era marks a definitive break
	with traditional interpretation, for good or for ill. Then again,
	others now seek to combine what can be known about worlds behind and
	before the Book of Genesis into a more comprehensive theological
	vision. In this course we will explore histories of Genesis in
	parallax, seeking above all to let them shape our understanding of
	the Book of Genesis as a core part of Jewish and Christian scripture
	from antiquity down through the present day.
}% end of course description
\objectives{% recall Bloom's taxonomy: http://www.celt.iastate.edu/teaching/effective-teaching-practices/revised-blooms-taxonomy
	By the end of the course students should be able to:
		remember and summarize events from the Genesis narratives in detail, locating them approximately by chapter;
		name and differentiate the most important ancestors (patriarchs and matriarchs) of Israel;
		explain different major models for the interpretation of Genesis, illustrating them with particular texts;
		use a variety of commentaries on Genesis from a number of different historical eras;
		distinguish text and interpretation by focusing on the afterlives of Adam and Eve;
		analyze claims made by different faith communities on the legacy of Abraham;
		critique specific arguments put forward by scholars about the meaning of select parts of Genesis;
		generate questions relevant to the interpretation of short passages from Genesis;
		construct an original argument in aid of the exposition of short passages from Genesis.
}% end of learning objectives

\ProvidesFile{preamble.tex}[2014/09/02 v1.1 -- Syllabus preamble]

% basic typography
\usepackage{fontspec,realscripts}
\defaultfontfeatures{Ligatures=TeX}
\setmainfont[%
  Numbers={Proportional,OldStyle},
  SmallCapsFeatures={RawFeature=+c2sc;+smcp,LetterSpace=1}]{Scala OT}
\newfontfamily\lining[Numbers={Lining}]{Scala OT}
\setsansfont[%
  SmallCapsFeatures={RawFeature=+c2sc;+smcp,LetterSpace=1}]{Scala Sans Pro}
\setmonofont[Scale=MatchLowercase]{Menlo}
\usepackage{sectsty}
\allsectionsfont{\sffamily\scshape}
\frenchspacing
\setlength{\emergencystretch}{3em} % prevent overfull lines

% custom font size and leading
\renewcommand\tiny{\fontsize{6}{9}\selectfont}
\renewcommand\scriptsize{\fontsize{7}{10}\selectfont}
\renewcommand\footnotesize{\fontsize{8}{10}\selectfont}
\renewcommand\small{\fontsize{9}{11}\selectfont}
\renewcommand\normalsize{\fontsize{10}{12}\selectfont}% base size
\renewcommand\large{\fontsize{11}{13}\selectfont}
\renewcommand\Large{\fontsize{12}{14}\selectfont}
\renewcommand\LARGE{\fontsize{16}{18}\selectfont}% "course syllabus \\ semester" benefits from more lead
\renewcommand\huge{\fontsize{20}{22}\selectfont}
\renewcommand\Huge{\fontsize{24}{26}\selectfont}

% layout packages: page, logo, tables
\usepackage[scale={0.6,0.8},
            xetex]{geometry}
\usepackage{graphicx}
\usepackage{array}     % allow insertions of column styling with >{}
\usepackage{booktabs}  % elegant horizontal rules in tables

% custom macros for a session count in the schedule of readings
\newcounter{session}
\newcounter{columns}
\newcounter{spancolumns}
\newcounter{midcolumns}
\newcounter{courseunit}
\newcommand\setcolumncount[2][0]{ % optionally set count to other than 0,
  \setcounter{session}{#1}        % e.g. to -1, or to a standing count
  \setcounter{columns}{#2}
  \setcounter{spancolumns}{#2-1}
  \setcounter{midcolumns}{#2-2}}
\newcommand\sessioncount{\stepcounter{session}\arabic{session}}
\newcommand\sessionskip[1]{\multicolumn{1}{@{}r@{ }}{#1}}
\newcommand\unit[1]{\multicolumn{\thecolumns}{c}{%
  \scshape\stepcounter{courseunit}\roman{courseunit}. #1}}
\newcommand\noclass[1]{\sessionskip{---.}&\multicolumn{\themidcolumns}{@{}l}{\itshape No Class: #1}}
\newcommand\rarr{\Large\char"EA31\hspace*{0.5pt}} % pointer finger w/ reminder tie in Scala OT
\newcommand\reminder[2]{\sessionskip{\rarr}&\multicolumn{\thespancolumns}{@{}l}{#1 \hfill #2}}

% color to match Tyndale's branding
\usepackage[usenames]{xcolor}
% predefined: black, white, red, green, blue, cyan, magenta, yellow
\definecolor{TyndaleURLs}{HTML}{0062A0} % links on tyndale.ca
\definecolor{TyndaleBlue}{cmyk}{1,1,0,.32}
\definecolor{TyndaleGold}{cmyk}{0,.27,1,0}
\definecolor{TyndaleRed}{cmyk}{0,1,.99,.04}
\definecolor{TyndaleBlack}{cmyk}{0,0,0,1}
\definecolor{TyndaleGreen}{cmyk}{.45,0,1,.24}
\definecolor{TyndaleOrange}{cmyk}{0,.79,1,0}
\definecolor{TyndaleAqua}{cmyk}{.47,0,.24,0}
\definecolor{TyndaleYellow}{cmyk}{.03,.03,.35,0}

% metadata (assumes a host of definitions are made in the main file)
\usepackage[setpagesize=false,     % leave this to geometry
            hyperfootnotes=false,  % fragile and distracting
            xetex]{hyperref}
\hypersetup{breaklinks=true,       % allow link text to break across lines
            colorlinks=true,       % colorlinks resets pdfborder to 0 0 0
            urlcolor=TyndaleURLs,  % for external links
            linkcolor=TyndaleRed,  % for normal internal links
            citecolor=TyndaleGold, % for bibliographical citations in text
            pdfauthor={\prof},
            pdftitle={\ccode: \ctitle},
            pdfsubject={Tyndale UC, \csemester},
            pdfcreator={github.com/danieldriver/syllabus}}
\urlstyle{same}                    % don't use monospace font for urls

% custom footlines
\usepackage{fancyhdr}
\pagestyle{fancy} % turn it on
\fancyhf{}        % reset everything
\renewcommand{\headrulewidth}{0pt} % remove header line as well
\lfoot{\sffamily\scshape\footnotesize\ctitle, \csemester}
\rfoot{\sffamily\scshape\footnotesize\prof\quad\thepage}

% language support: polyglossia needs to go last
\usepackage{polyglossia}
\setdefaultlanguage{english}
\setotherlanguage{hebrew}
\newfontfamily\hebrewfont[Script=Hebrew]{SBL BibLit}

% protect positioning of margin table in policy/grades -- needs to come after polyglossia
\usepackage{marginfix}

% gratuitous with custom title page, but useful as a fallback
\title{\ccode: \ctitle}
\author{\professor}
\date{\semester}


\begin{document}
\ProvidesFile{title.tex}[2013/09/06 v1.0 -- Syllabus title page]

\begin{titlepage}
  \begin{center}

    \begin{minipage}{\textwidth}
      \parbox[t]{0.5\textwidth}{
        \mbox{}\\[-\baselineskip] % dummy line to align parboxes
        \includegraphics[width=0.5\textwidth]{.syllabus/includes/TyndaleUC}}
      \hfill
      \parbox[t]{0.4\textwidth}{
        \raggedleft\LARGE\sffamily Course Syllabus\\
        \semester}
    \end{minipage}

    \vfill

    {\sffamily\textsc{\LARGE\MakeLowercase\coursecode}\\[2ex]
      \bfseries\itshape\Huge\coursetitle}

    \vfill

    \begin{tabular}{>{\bfseries}rl>{\bfseries}rl}
      \toprule
      Instructor & \professor                   & Course  & \liningnums{\coursecode}~\texttt{(v\version)} \\
      \midrule
      Phone      & \texttt{\phone}              & Meets   & \meeton                                       \\
      Email      & \href{mailto:\email}{\email} & Time    & \meetat                                       \\
      Office     & \officenumber                & Room    & \meetin                                       \\
      Hours      & \officehours                 & Website & \href{http://\website}{\website}              \\
      \bottomrule
    \end{tabular}

    \vfill

    \begin{description}
      \item[Commuter Hotline]
        Class cancellations due to inclement weather or illness will
        be announced on the commuter hotline at \texttt{416.226.6620
        x2187}. Alternately, weather cancellation information is posted
        at \href{http://tyndale.ca/weather}{tyndale.ca/weather}.
      \item[MyTyndale.ca]
        This course may have materials stored on its website, such as
        handouts or readings that may be needed in order to complete
        assignments. Students are responsible for checking these course
        pages on a regular basis. Here, too, students are able to view
        their grades throughout the semester. For more information see
        Section~\ref{mytyndale}, below.
      \item[Mail]
        Students are responsible for information communicated through
        their campus mailboxes and student e-mail accounts. A mailbox
        directory hangs beside the mailboxes. For more information
        contact the Registrar's office.
    \end{description}

  \end{center}

  \section{Course Description}
  \label{description}

  \emph{From the Academic Calendar:} \cdescription

\end{titlepage}
\setcounter{page}{2} % count the title page as page 1


\section{Learning Objectives}
\label{objectives}
\edobject

\section{Required Texts \& Materials}
\label{texts}

The following texts are required. Students are strongly encouraged to
purchase their own copies. Library copies that are not reference works
will be placed on a 2-hour reserve.

\begingroup
\renewcommand{\section}[2]{}% temporarily remove the section heading
\begin{thebibliography}{Anderson}% use the longest item in the bibliography

	\bibitem[NRSV]{nrsv} M.\,D. Coogan, ed.
    \emph{New Oxford Annotated Bible with Apocrypha: NRSV}. 4th ed.
    Oxford / New York: Oxford University Press, 2010.
    \textsc{isbn} 978-0195289602.

%	\bibitem[NJPS]{njps}
%	Adele Berlin and Marc Zvi Brettler, eds.
%	\emph{The Jewish Study Bible: Second Edition}.
%	Oxford / New York: Oxford University Press, 2014.
%	\textsc{isbn} 978-0199978465.

	\bibitem[Anderson]{anderson} Gary A. Anderson.
	\emph{The Genesis of Perfection: Adam and Eve in Jewish and Christian Imagination}.
	Louisville: Westminster John Knox, 2002.
	\textsc{isbn} 978-0664226992.

	\bibitem[Levenson]{levenson} Jon D. Levenson.
	\emph{Inheriting Abraham: The Legacy of the Patriarch in Judaism, Christianity, and Islam}.
	Princeton: Princeton University Press, 2012.
	\textsc{isbn} 978-0691163550.

	\bibitem[Provan]{provan} Iain Provan.
	\emph{Discovering Genesis: Content, Interpretation, Reception}.
	Grand Rapids: Eerdmans, 2016.
	\textsc{isbn} 978-0802872371.

%	\bibitem[tk]{tk} tk.
%	\emph{tk}.
%	tk: tk, 20tk.
%	\textsc{isbn} 978-tk.

\end{thebibliography}
\endgroup

An acceptable alternative study Bible is the NJPS: Adele Berlin and Marc
Zvi Brettler, eds., \emph{The Jewish Study Bible: Second Edition}
(Oxford / New York: Oxford University Press, 2014). A reference copy is
available in the library, and it is well worth consulting.

\section{Supplementary Texts}
\label{supplementary}

Supplementary readings will be recommended throughout the semester.
Excerpts from this literature may either be placed on reserve or made
available for download through the course website. Outstanding students
will also pursue leads on their own, making maximum use of notes and
bibliographies in the required textbooks.

\begin{multicols}{3}[Major modern commentaries on Genesis include:]%
\footnotesize\noindent
M.\,M. Kalisch (1858)\\
J.\,P. Lange (ET 1868)\\
F. Delitzsch (\textsuperscript{5}1899)\\
A. Dillmann (\textsuperscript{6}1892)\\
H. Holzinger (1898)\\
S.\,R. Driver (1904)\\
J. Skinner (1910)\\
H. Gunkel (\textsuperscript{4}1917)\\
O. Procksch (\textsuperscript{3}1924)\\
B. Jacob (1934)\\
A. Richardson (1953)\\
A. Clamer (1953)\\
U. Cassuto (1961--64)\\
E.\,A. Speiser (1964)\\
W. Zimmerli (\textsuperscript{3}1967, 76)\\
D. Kidner (1967)\\
G. von Rad (ET \textsuperscript{2}1972)\\
R. Davidson (1973)\\
C. Westermann (1974)\\
B. Vawter (1977)\\
N. Sarna (1989)\\
J. Ebach [37--50, HthKAT] (2007)\\
T. Longman (2016)\\
C. Dohmen [1--11, KB] (2017)
\end{multicols}

\begin{multicols}{2}[Commentaries on Genesis in the AST library include:]%
\footnotesize\noindent
Targums: {\scshape BS 709.2 B5 1987 vv. 1A, 1B, 6} (c. 150--?)\\
Didymus the Blind: {\scshape BS 1235 D4913 2016} (d. 398)\\
Augustine of Hippo: {\scshape BS 1235 A8413 1991} (d. 430)\\
Bede, the Venerable: {\scshape BS 1235 B43 2008} (d. 735)\\
ACCS, Genesis 1--11: {\scshape BS 1235.3 G46 2001}\\
ACCS, Genesis 12--50: {\scshape BS 1235.53 G46 2002}\\
RCS, Genesis 1--11: {\scshape BS 1235.53 G455 2012}\\
Luther, Martin: {\scshape BR 330 E5 1955 vv.1--8} (d. 1546)\\
Calvin, John: {\scshape BS 1235 C293 1948} (d. 1564)\\
Patrick, Simon: {\scshape BS 1235 P36 1695}\\
Henry, Matthew: {\scshape BS 490 H4 1961} (d. 1714)\\
Hershon, Paul Isaac: {\scshape BS 1235 H47} (1883)\\
Delitzsch, Franz: {\scshape BS 1235 D4} (1888)\\
Skinner, John: {\scshape BS 1235 S45 1917, 1963} (1930)\\
Richardson, Alan: {\scshape BS 1235 R66 1959}\\
Cassuto, Umberto: {\scshape BS 1235.3 C3} (1961)\\
Herbert, Arthur: {\scshape BS 1235.3 H4} (1962)\\
Speiser, E. A.: {\scshape BS 1233 S64 1964}\\
Kidner, Derek: {\scshape BS 1235.3 K47} (1967)\\
von Rad, Gerhard: {\scshape BS 1235.3 R3213 1972}\\
Plaut, Gunther: {\scshape BS 1225.3 P55 v. 1} (1974)\\
Vawter, Bruce: {\scshape BS 1235.3 V38 1977}\\
Davidson, Robert: {\scshape BS 1235.3 D3 1973, 1979}\\
Leibowitz, Nehama: {\scshape BS 1235.3 L413 1981}\\
Wenham, Gordon : {\scshape BS 491.2 W67 1982}\\
Brueggemann, Walter: {\scshape BS 1235.3 B78} (1982)\\
Maher, Michael: {\scshape BS 1235.3 M346 1982}\\
Westermann, Claus: {\scshape BS 1235.3 W43213 1986, 1987, 1994}\\
Sarna, Nahum: {\scshape BS 1235.3 S325 1989}\\
Hamilton, Victor: {\scshape BS 1235.3 H32 1990, 1995}\\
Scullion, John: {\scshape BS 1235.3 S37 1992}\\
Fox, Everett: {\scshape BS 1223 A3 F68 1995}\\
Hamilton, Victor: {\scshape BS 1235.3 H323 1995}\\
Hartley, John: {\scshape BS 1235.3 H37 2000}\\
Towner, Sibley: {\scshape BS 1235.3 T69 2001}\\
Walton, John: {\scshape BS 1235.53 W35 2001}\\
Cotter, David: {\scshape BS 1235.52 C68 2003}\\
Briscoe, Stuart: {\scshape BS 1151.2 C66 2004}\\
McKeown, James: {\scshape BS 1235.53 M35 2008}\\
Arnold, Bill: {\scshape BS 1235.53 A76 2009}\\
Goldingay, John: {\scshape BS 1235.53 G65 2010}\\
Reno, Russell: {\scshape BS 1235.53 R46 2010}\\
Cook, Joan: {\scshape BS 1235.53 C66 2011}\\
De La Torre, Miguel: {\scshape BS 1235.53 D4 2011}\\
Coleson, Joseph: {\scshape BS 1235.53 C64 2012}\\
Longman, Tremper: {\scshape BS 1235.53 L66 2016}
\end{multicols}

\section{Course Outline}
\label{outline}

We will adhere to the schedule in \autoref{schedule} as closely as
possible, though the professor reserves the right to adjust it to suit
the needs of the class.

\setcolumncount{7}% set up \sessioncount, \unit{}, \noclass{}, and \reminder{memo}{date} macros
\begin{table}[htbp]% set to `p' to put the schedule on its own page
  \centering
  \begin{tabular}{>{\sessioncount.}r@{ }lccccr}% make sure the column config agrees with \setcolumncount
	\toprule
	\sessionskip{\textbf{\S}.}&\textbf{Primary}&\textbf{Anderson}&\textbf{Levenson}&\textbf{Provan}&\textbf{Seminar}&\textbf{Date}\\
	\midrule

		& Genesis 49     &          &       &      &    &  6 Jan. \\
		& Genesis 9      & Intro--1 &       & 1--3 & DD & 13 Jan. \\
		& Genesis 1      & 2--3     &       & 4    & MT & 19 Jan. \\
		& Genesis 2      & 4--5     &       & 5    & AM & 27 Jan. \\
		& Genesis 3      & 6--7     &       & 6    & MT &  3 Feb. \\
		& Genesis 4--5   & 8--End   &       & 7    & AM & 10 Feb. \\
		& Genesis 6--11  &          & Intro & 8    & MT & 17 Feb. \\
	\noclass{Term Break (Monday to Friday)}             & 24 Feb. \\
		& Genesis 12--15 &          & 1--2  &      & AM &  3 Mar. \\
		& Genesis 16--25 &          & 3--4  & 9    & MT & 10 Mar. \\
		& Genesis 26--36 &          & 5     & 10   & AM & 17 Mar. \\
		& Genesis 37--45 &          & 6     & 11   & MT & 24 Mar. \\

	\reminder{Exegetical paper is \textbf{due} at the start of class twelve}{} \\

		& Genesis 46--50 &          &       &      & AM & 31 Mar. \\

	\reminder{End of Term: Final marks are due for all courses}{12 Apr.} \\

	\bottomrule
  \end{tabular}
  \caption{Schedule of Readings}
  \label{schedule}
\end{table}

See the AST website for a list of other \href{http://www.astheology.ns.ca/students/academic-dates.html}{important dates}.

\section{Evaluation and Grade Structure}
\label{evaluation}

\begin{enumerate}

	\item The professor will set \textbf{reading prompts} throughout the
	semester, as often as once a week. Prompts will consist of questions
	drawn from the secondary reading, and so presuppose that students
	have read the assigned material carefully. Points may be deducted if
	it becomes evident that the reading has not been taken seriously.

	\item Every other week, students will lead \textbf{seminars} on
	single chapters from the Book of Genesis. They should take the lead
	by briefly outlining the chapter in question, situating it relative
	to the rest of Genesis, and reporting on their research into part of
	that chapter. Then, in the balance of the time, they should helping
	the class elaborate, reflect on, test, challenge, or extend a
	central idea or two. Students are encouraged to align their seminar
	preparations with their choice of texts for their exegetical paper.

	\item A major \textbf{exegetical paper} will facilitate an extended
	exercise in biblical interpretation. It should focus on a short
	passage from Genesis, should interact with the secondary readings
	and further research into the passage in question, and should be
	3--4,000 words long. It is due at the start of the last class.

%	\begin{enumerate}
%
%		\item The \textbf{first paper} will focus on a passage from
%		Genesis 1--11 and should interact carefully with
%		\cite{anderson}, among others.
%
%		\item The \textbf{second paper} will focus on a passage from
%		Genesis 12--50 and should interact carefully with
%		\cite{levenson}, among others.
%
%	\end{enumerate}

\end{enumerate}

The breakdown for the semester's total work is shown in
\autoref{grade-dist}.

\begin{table}[htbp]
  \centering
  {\lining
  \begin{tabular}{lr}
    \toprule
    Reading Prompts  & 30\% \\
    Student Seminars & 30\% \\
    Exegetical Paper & 40\% \\
    \bottomrule
  \end{tabular}}
  \caption{Distribution of Grades}
  \label{grade-dist}
\end{table}

\ProvidesFile{grades.tex}[2023/09/06 v2.5 -- Course policy]

AST's \href{https://www.astheology.ns.ca/students/resources.html}{Academic
Calendar} provides guidelines and criteria for academic assessment. Marks are
assigned by letter grade using these benchmarks. Note that graduate and
undergraduate programs are currently assessed on a different scales at AST (see
\autoref{grade-syst}). For graduate students, a grade of C is deemed
``Unsatisfactory,'' whereas for undergraduate (BTh) students, a C is
``Satisfactory'' and the grade of D is possible for ``Marginal'' performance.
Other grades have different definitions, too.

\begin{table}[htbp]
 \centering
 {\lining\footnotesize
 \begin{tabular}{lll|ll}
   \toprule
               & \multicolumn{2}{c|}{Graduate Courses}   & \multicolumn{2}{c}{Undergraduate Courses} \\
               \cline{2-5}
   A+          & 94--100    & Exceptional               & 90--100    & Excellent                 \\
   A           & 87--93     & Outstanding               & 85--89     & Excellent                 \\
   A\char"2212 & 80--86     & Excellent                 & 80--84     & Excellent                 \\ [1ex]
   B+          & 77--79     & Good                      & 75--79     & Good                      \\
   B           & 73--76     & Acceptable                & 70--74     & Good                      \\
   B\char"2212 & 70--72     & Marginally Acceptable     &            &                           \\ [1ex]
   C           & 60--69     & Unsatisfactory            & 60--69     & Satisfactory              \\
   D           &            &                           & 50--59     & Marginal                  \\
   F           & 0--59      & Failure                   & 0--49      & Failure                   \\
   FP          & 0          & Failure due to Plagiarism & 0          & Failure due to Plagiarism \\
   \bottomrule
 \end{tabular}}
 \caption{Outline of Grades at AST}
 \label{grade-syst}
\end{table}

\ProvidesFile{other.tex}[2018/0/03 v2.4 -- Course policy]

\section{Other Course Policy}
\label{policy}

Late work will not be accepted, except in genuinely extenuating
circumstances. Students must submit something before the deadline if
they wish to receive credit. Unless I state otherwise, written
assignments are to be uploaded by 11:59 \PM\ on the date indicated.

Essay submissions must be typewritten, double-spaced, and formatted as
PDFs. They should be free from error. In this course they should follow
either \emph{The SBL Handbook of Style} (refer to the online
\href{https://www.sbl-site.org/assets/pdfs/pubs/SBLHSsupp2015-02.pdf}{
Student Supplement}), or \emph{The Chicago Manual of Style}, per AST's
policy for major papers. As a reminder, AST also upholds an Inclusive
Language Policy.

\href{http://www.eerdmans.com/Pages/Item/59043/Commentary-Statement.aspx}{Plagiarism},
if \href{https://www.theguardian.com/world/2013/feb/09/german-education-minister-quits-phd-plagiarism}{detected},
will result in failure of the course.

Students should request permission to record a class or lecture. If
permission is granted, or if recordings are provided (as in the case of
an online course), I stipulate that all recordings be for personal use
only. They may not be shared or distributed.

If you have abilities or disabilities that require modifications to the
assessment process or other aspects of this course, please advise the
course instructor as soon as possible.

Finally, I encourage the conscientious use of laptops, tablets, and
other technology in my classes. In classroom settings, realize that, as
\href{http://dx.doi.org/10.1016/j.compedu.2012.10.003}{cognitive
psychologists have demonstrated}, ``laptop multitasking hinders
classroom learning for both users and nearby peers.'' Do your part to
foster an environment of open dialogue by honouring the presence of your
classmates. In online settings, consider both the physical environment
in which you choose to work and the virtual environment that you help
create through your participation in various forums. Let your engagement
in this course be marked by rigour and charity alike.

\end{document}